% RESEARCH_DRAFT_SHA256: 3b77498e85b930d7f93956006cc9fbfb19d750fbb66ad870cfafa62a1effa97d
% Layout and source-reference wrappers expanded for the web renderer.
\documentclass[11pt]{article}
\usepackage[margin=25mm]{geometry}
\usepackage{fontspec,unicode-math}
\setmainfont{Latin Modern Roman}
\setmathfont{Latin Modern Math}
\usepackage{amsmath,amssymb,mathtools,xurl,hyperref}
\hypersetup{colorlinks=true,urlcolor=blue,linkcolor=blue}
\setcounter{secnumdepth}{0}
\setlength{\parindent}{0pt}
\setlength{\parskip}{5pt}
\setlength{\emergencystretch}{3em}
\newcommand{\sref}[2]{\hyperlink{source:#1:#2}{[S#2]}}
\newcommand{\cataloguescope}[1]{\paragraph{Source record.}#1}
\newcommand{\floor}[1]{\left\lfloor#1\right\rfloor}
\begin{document}
\section*{Research-attempt provenance}
This standalone research attempt retains the complete original definition body below,
including its historical source and preparation notes. Those notes describe the
earlier catalogue preparation; the new attempt and its verification are identified
separately after that body. No original ID, catalogue status, or source statement is changed.

% UnsolvedMath ID 20002161; source catalogue code AIM-INFRASTRUCTURE-0061
\section{Coherent numbering for extracted PreTeXt content}\label{20002161}
{\small UnsolvedMath ID 20002161 · Computer Science}
\subsection{Definitions and mathematical statement}

\paragraph{Shared notation.}Standard number systems are $\mathbb N,\mathbb Z,\mathbb Q,\mathbb R,\mathbb C$. Specialized notation and index ranges follow the cited source and the definitions below.


Displayed numbering assigns labels such as chapter, theorem, or figure numbers according to a hierarchy. Preserving numbers retains labels from the original document; renumbering computes labels in the extracted document. Structural identity and displayed number are different properties of an item.
\paragraph{Record scope and qualifications.}
The catalogue title is retained as metadata. Named certificates, contracts, bounds or frameworks in that title are not asserted as results of the original source extract.
The numbering hierarchy, identifier mapping, and effects on cross-references need an implementation specification.

Catalogue screening limitation: PreTeXt software design, documentation, or publication feature request; no mathematical problem is specified by the supplied passage.

The source snapshot labels this record partially solved; this is catalogue metadata rather than independent verification of a mathematical result or the current completion of a project.
\paragraph{Statement recorded in the supplied source.}
Numbering when extracting. Preserve or renumber? Levels of numbering (figure for example). Both need to be publisher options. \href{https://github.com/PreTeXtBook/community/wiki/Moderated-Problem-Session}{[S2]}
\paragraph{Source status.}
The supplied catalogue labels this source record partially solved; that metadata was checked against the supplied snapshot on 2026-10-11. The label is preserved for provenance and is not an independent verification of present open status or mathematical correctness.
\subsection{Short English statement}
Provide publication options for preserving or recomputing numbering when extracting content.
\subsection{Sources}
\begin{itemize}
\item[S1] Ulam AI and the UnsolvedMath Contributors, UnsolvedMath: A Curated Collection of Open Mathematics Problems, record 20002161 (AIM-INFRASTRUCTURE-0061). Catalogue curation and metadata: CC BY 4.0. The source extract is retained; definitions, English summaries and qualifications are editorial additions.
\url{https://huggingface.co/datasets/ulamai/UnsolvedMath}.
\item[S2] AIM Workshop Problem Lists, source linked by catalogue record AIM-INFRASTRUCTURE-0061.
\url{https://github.com/PreTeXtBook/community/wiki/Moderated-Problem-Session}.
\end{itemize}
\paragraph{Preparation and review.}
The original source and UnsolvedMath contributors retain attribution for the source extract. Codex (GPT-6; exact serving revision not exposed) assisted with definitions, English summaries, scope notes and formatting. No new human mathematical authorship or independently verified solution is asserted. Permanent source ID 20002161 and code AIM-INFRASTRUCTURE-0061 checked against the complete supplied registry and current repository. Distinct upstream records, including repeated questions and broad or damaged extracts, are retained for complete coverage of the source CS category. A repeated record is not a new independent conjecture. Complete mathematical-equivalence and current-literature review has not been performed. Preparation and final review follow contribution prompts C and D, with source gaps disclosed. Structural and rendering checks do not establish a new theorem or open status.
\label{end:20002161}

\section{Literature and context check}
The primary community-session entry was read. It is an implementation or documentation agenda, and its editorial catalogue title supplies no theorem specification. The following explicitly chosen model tests one route; no actual PreTeXt conversion, documentation release, or user study is claimed.
\section{Research attempt: Extraction numbering and the impossibility of one universal display number}
Take a full publication whose figures $a,b,c$ have display numbers $1,2,3$. Extract only $c$. A preserve-number policy assigns $c\mapsto3$; a compact-renumber policy assigns $c\mapsto1$. A system storing only one unqualified number for $c$ cannot satisfy both policies at once, since $1\ne3$. Thus the semantic identifier and the display number cannot be interchangeable if both publisher options are required.

For a finite extracted ordered set $S$ of identifiers, define $N_{\rm keep}(i)=N_{\rm full}(i)$ and $N_{\rm compact}(i)=1+|\{j\in S:j\text{ precedes }i\}|$. The second is injective: if $i$ precedes $j$, the predecessor set of $j$ contains that of $i$ and also $i$, so its count is strictly larger. The first is injective if the full numbering is injective within the relevant counter scope. Both maps are computable and give coherent references when every heading and reference in that rendition reads the same selected map at the semantic identifier.

Numbering levels add scope. If each chapter restarts figure numbering at one, local number 1 occurs in several chapters; it identifies a figure only together with its chapter scope, or as a qualified display such as $(\text{chapter},1)$. Combining two separate books can still duplicate that pair, so a publication identity may also be needed for cross-publication references. These collisions are exact examples, not reasons to prohibit local counters. They show that display labels must be interpreted under an explicit numbering context.

The maps do not decide whether preserving or renumbering is better. They prove that either can be consistent in a selected context and that mixing their numbers in one reference lookup can select the wrong object.
\section{Verification and adversarial audit}
The finite examples and implications are checked directly in the argument. No production PreTeXt build or conformance test was run. The proof applies only to the declared model, with its input and identity assumptions.
\paragraph{Audit corrections.} Injectivity is stated within an explicit counter scope. A policy-dependent display label is not used as the permanent semantic identifier.
\paragraph{External computation.} None; symbolic arguments and boundary checks only.
\section{Scope and remaining gaps}
The result is an auxiliary route analysis, not completion of the original software or documentation agenda. Real versions, supported syntax, transformations, and usability requirements need concrete implementation and testing. No mathematical novelty or current feature-completion claim is made.
\paragraph{Integration review.} Codex (GPT-6; exact serving revision not exposed) reviewed the written arguments and expanded layout and source-reference wrappers. Repository metadata and the production web renderer were checked using Python 3.13 and JavaScript with MathJax 3.2.2. This model-assisted review is not independent human proof verification; no PDF compilation or formal proof checking is asserted.

\section{Completion Estimate}
\noindent\textbf{COMPLETION: 0\%}
The calculations below the preserved source statement investigate a concrete route,
but establish no substantiated new progress toward the entire original target.
Known special cases and failed sufficient conditions are not credited as a solution.
\section{References}
The original source references and catalogue attribution are preserved in the
Sources section above. The following URLs were actually consulted for this attempt;
retention of an additional original reference does not claim fresh consultation of it.
\begin{itemize}
\item \url{https://github.com/PreTeXtBook/community/wiki/Moderated-Problem-Session} (consulted 2026-10-11). 
\end{itemize}
\section{Model and contributor disclosure}
OpenAI gpt-6-astra, requested reasoning configuration ultra, generated this research
attempt on 2026-10-11. The exact serving revision is not exposed. Codex GPT-6
performed orchestration and repository-format review; its exact serving revision
is also not exposed. The mathematical route, boundary checks and final adversarial
self-audit followed contribution prompts A then B in
\texttt{docs/contribution\_prompts.md}, after reading \texttt{CONTRIBUTING.md}.
Python 3.14.2 assembled the standalone documents; any mathematical computation is
specified in Verification. No theorem prover was used. No human mathematical
authorship, substantial human proof correction, or independent human proof checking
is asserted. Model-assisted review, source retrieval, and structural checks do not
establish correctness or novelty. Full private reasoning transcripts are unavailable
as reproducibility artifacts; the complete arguments, inputs and outcomes actually
used are stated here or in the companion batch audit.
\section{Final status}
Unresolved as a new research contribution. This is a concrete route investigation,
not a claim of new complete resolution or a current-status certification. A known
answer reproduced for context does not become a new contribution. The catalogue's
source-status label is preserved; this attempt status does not assert that a
source-labelled solved question is open.
\end{document}
